% =============================================================================
% APPENDIX D: FAST External Validation Methodology
% =============================================================================

\clearpage
\section{FAST External Validation Methodology}
\label{app:fast-validation}

Phase~7 of the experiment validates model-ladder predictions against
FEMA FAST (Flood Assessment Structure Tool) engineering damage estimates.
FAST applies Hazus depth--damage curves to individual structures from the
USACE National Structure Inventory~2.0 (NSI), producing per-structure
economic loss estimates that are aggregated to ZCTA level.  This appendix
documents the event-level matching procedure, which resolves a
granularity mismatch between model predictions and synthetic hazard
scenarios.

\subsection{Granularity Mismatch}
\label{app:fast-mismatch}

The model ladder produces predictions at the \texttt{(zcta\_id, event)}
grain: each ZCTA appears once per historical flood event (e.g., Harvey~2017,
Beryl~2024).  FAST produces loss estimates at the \texttt{(zcta\_id,
return\_period)} grain: one estimate per ZCTA per synthetic hazard
scenario (e.g., 100-year rainfall, Category~3 surge).

A naive approach would collapse events by averaging all predictions per
ZCTA and correlating against a single return period.  This conflates
events of very different intensities (Harvey is a $>$1000-year rainfall
event; Beryl is $\sim$25-year) into one number, producing a semantically
invalid correlation.

\subsection{Event-Level Matching Protocol}
\label{app:fast-matching}

We resolve this mismatch by mapping each historical event to its closest
synthetic return period \emph{before} computing correlations.  The
mapping is fixed prior to Phase~7b execution (no data-adaptive
selection).

\paragraph{Pluvial scenarios (Houston, NYC).}
FloodSimBench provides physics-based 2D shallow-water flood depth
rasters at 10 design rainfall return periods (1--1000\,yr), defined by
6-hour rainfall totals.  We match each historical event by comparing
its observed NOAA MRMS Stage~IV peak 6-hour rainfall total against the
FloodSimBench bins:

\begin{table}[h]
\centering
\small
\caption{Event-to-return-period assignments (pluvial scenarios).
  Rainfall values are peak 6-hour MRMS totals; FloodSimBench bins
  are the nearest 6-hour design rainfall.}
\label{tab:event-rp-pluvial}
\begin{tabular}{llrrl}
\toprule
Scenario & Event & Peak 6h (mm) & Matched RP & FloodSimBench bin \\
\midrule
Houston & Harvey 2017    & $\sim$180 & 1000-yr & 181\,mm/6h \\
Houston & Imelda 2019    & $\sim$140 & 200-yr  & 138\,mm/6h \\
Houston & Beryl 2024     & $\sim$100 & 25-yr   & 98\,mm/6h  \\
NYC     & Ida 2021       & $\sim$120 & 100-yr  & 123\,mm/6h \\
NYC     & Henri 2021     & $\sim$45  & 1-yr    & 48\,mm/6h  \\
\bottomrule
\end{tabular}
\end{table}

\paragraph{Surge scenario (Southwest Florida).}
NOAA SLOSH Maximum of MEOWs (MOM) provides category-specific surge
inundation rasters.  Each event is matched to its Saffir-Simpson
category at landfall:

\begin{table}[h]
\centering
\small
\caption{Event-to-return-period assignments (surge scenario).
  Category assignments follow NHC best-track data (HURDAT2).}
\label{tab:event-rp-surge}
\begin{tabular}{llll}
\toprule
Scenario & Event & Landfall Cat. & SLOSH MOM \\
\midrule
SW Florida & Ian 2022    & Category 4 & Cat-4 HIGH tide \\
SW Florida & Helene 2024 & Category 4 & Cat-4 HIGH tide \\
SW Florida & Milton 2024 & Category 3 & Cat-3 HIGH tide \\
\bottomrule
\end{tabular}
\end{table}

\subsection{Correlation Protocol}
\label{app:fast-protocol}

For each (scenario, event, matched return period) triple:
\begin{enumerate}
  \item Load FAST ZCTA aggregates (\texttt{fast\_total\_loss\_zcta})
        for the matched return period.
  \item Load model predictions for the specific event at each
        representation level (R0, R1, R2), averaging across CV folds
        to produce one prediction per ZCTA.
  \item Compute Spearman $\rho$(prediction, FAST loss) per level.
  \item Compute Spearman $\rho$(observed NFIP claims, FAST loss) as
        a ceiling---the maximum correlation achievable if the model
        perfectly predicted claims.
\end{enumerate}

\noindent
This yields 8 event-level correlations (3 Houston + 2 NYC + 3 SW~Florida)
per representation level, plus 8 ceiling values.

\subsection{Robustness}
\label{app:fast-robustness}

We assess robustness along two axes:
\begin{itemize}
  \item \textbf{Cross-event consistency}: Does the level ranking
        (R0 vs R1 vs R2) persist across events within a scenario?
  \item \textbf{Cross-scenario consistency}: Does the level ranking
        persist across scenarios with different flood mechanisms
        (pluvial vs surge)?
\end{itemize}

\noindent
For transparency, we also report the naive collapsed correlation
(all events averaged per ZCTA) so the reader can assess the impact
of event-level matching on reported effect sizes.

\subsection{Limitations}
\label{app:fast-limitations}

\begin{itemize}
  \item \textbf{Approximate matching.}  The event-to-return-period
        mapping uses peak 6-hour MRMS rainfall as a proxy.  Actual
        flood depth depends on antecedent soil moisture, drainage
        capacity, and storm duration---all of which differ between
        the synthetic and historical scenarios.
  \item \textbf{Partial coverage.}  FloodSimBench covers Houston
        (7~tiles, full metro) and NYC (2~tiles, Manhattan only).
        SLOSH MOM is coastal-only.  ZCTAs outside raster coverage
        receive no FAST estimate and are excluded from correlation.
  \item \textbf{No New Orleans or Riverside.}  Neither scenario has
        synthetic hazard rasters available; they are excluded from
        FAST validation entirely.
  \item \textbf{SLOSH is category-based, not ARI-based.}  The
        Saffir-Simpson-to-SLOSH mapping is deterministic, not
        probabilistic.  Two Category~4 storms produce the same SLOSH
        MOM estimate despite different surge footprints.
\end{itemize}
